% ---- pg-529 （书页 517） ----

\begin{mvtable}
\mvtabcap{续表}
\begin{longtable}{ccccc}\toprule
$x_1$ & $x_2$ & $x_3$ & $x_4$ & $x_5$\\
\midrule
\endfirsthead
\multicolumn{5}{c}{\zihao{6}续表}\\
\toprule
$x_1$ & $x_2$ & $x_3$ & $x_4$ & $x_5$\\
\midrule
\endhead
\bottomrule
\endlastfoot

8.0 & 14.0 & 0.30 & 4.32 & 7.87\\
8.4 & 18.0 & 0.20 & 4.38 & 7.98\\
10.0 & 18.0 & 0.10 & 3.06 & 7.67\\
7.3 & 15.0 & 0.05 & 3.76 & 6.84\\
9.5 & 22.0 & 0.30 & 3.98 & 5.02\\
8.4 & 15.0 & 0.20 & 5.02 & 10.12\\
8.4 & 17.0 & 0.20 & 4.42 & 8.25\\
9.5 & 25.0 & 0.50 & 4.44 & 5.95\\
7.2 & 22.0 & 1.00 & 4.70 & 3.49\\
4.0 & 12.0 & 0.50 & 5.71 & 6.32\\
6.7 & 52.0 & 0.50 & 4.80 & 3.20\\
9.0 & 27.0 & 0.30 & 3.69 & 3.30\\
7.8 & 29.0 & 1.50 & 6.72 & 5.75\\
4.5 & 41.0 & 0.50 & 3.33 & 2.27\\
6.2 & 34.0 & 0.70 & 7.56 & 6.93\\
5.6 & 20.0 & 0.50 & 5.07 & 6.70\\
9.0 & 17.0 & 0.20 & 4.39 & 8.33\\
8.4 & 20.0 & 0.10 & 3.74 & 3.77\\
9.5 & 19.0 & 0.50 & 3.72 & 7.37\\
9.0 & 20.0 & 0.50 & 5.97 & 11.17\\
6.2 & 16.0 & 0.05 & 4.23 & 4.18\\
7.3 & 20.0 & 0.50 & 4.39 & 3.50\\
3.6 & 15.0 & 0.70 & 7.00 & 4.82\\
6.2 & 34.0 & 0.07 & 4.84 & 2.37\\
7.3 & 22.0 & 0.00 & 4.13 & 2.70\\
4.1 & 29.0 & 0.70 & 5.78 & 7.76\\
5.4 & 29.0 & 0.20 & 4.64 & 2.65\\
5.0 & 34.0 & 0.70 & 4.21 & 6.50\\
6.2 & 27.0 & 0.30 & 3.97 & 2.97\\
\end{longtable}
\end{mvtable}

\noindent 11.31\quad 参考\hyperref[tab:11.2]{表 11.2} 中的鲑鱼数据.

(a) 对两组鲑鱼的二元数据作图. 这些数据点的大小和方向是否大致相同? 对阿拉斯加和加拿大鲑鱼数据, 二元正态分布与等协方差矩阵是否是可以接受的总体模型?

(b) 对于等先验概率和等代价条件下的两个正态总体, 用线性判别函数构造判别得分的点图. 年轮直径指标是否将两个总体分离得很好? 请解释.

(c) 分别对雄性和雌性鱼重复例 11.7 中的分析. 判别两类雄鱼是否比判别两类雌鱼容易些? 性别是否可以作为的一个有用的判别量?

\noindent 11.32\quad \hyperref[tab:11.8]{表 11.8} 所列数据与例 11.3 中所用的 A 型血友病基因携带者数据类似(见文献[15]). 利用这些数据:

(a) 对两个组检查二元正态性假定.

(b) 假定等先验概率, 求样本线性判别函数, 并用提留方法估计失误率.
